\section{模型的评价与总结}
\label{sec:evaluation}

\subsection{模型的有效性与优点}
三个模型以同一张图的整图单Task单核工期为速度比基准，再逐步改变张量的可复用范围：问题一统计Task接收域，问题二统计核心接收域，问题三在核心驻留之上加入Cache查询时序。100张图的统一基准均有官方评价结果；主矩阵覆盖100张图、1--5核和A/B/C三种场景，共1500个图--核数--场景组合，其中A的单核100组用于基准核验，正式多核组合为1400组。
每个组合保留一条最终官方结果；候选成功率以候选记录数为分母，组合覆盖率以1500个图--核数--场景组合为分母。本文的工期、搬运和命中字段均取相应组合的官方结果。

对图$g$、核数$k$、场景$r$，真实方案速度比为$S_{gkr}=T_{g0}/T_{gkr}$。题目要求A、B单核曲线点显示为1，CSV仍保留实际方案工期；B单核真实速度比的100图均值为0.9878。A/B的1000个存储结果有180个速度比低于1；正式A/B的900个方案结果中为177个，若只看2至5核则为126/800。图上统一的单核起点方便对比，逐图工期和慢例分析使用真实数值。

\subsection{基线与受控消融对照}
先以同一整图单Task单核工期作为统一基线，再固定图、核数和计划逐项改变一个因素。表\ref{tab:baseline-ablation}将相对A0的方案速度比与B/C同计划配对比值分列；$B(X_B^*)\to C(X_B^*)$固定计划只改变L2配置，$C(X_B^*)\to C(X_C^*)$固定C配置只改变切图与核内顺序，二者构成受控消融。首个成功候选到最终候选的比较另记为搜索增益，不承担单一机制的消融结论。
\begin{table}[H]\centering\small
\caption{五核100图的统一基线与受控消融（两种配对口径）}
\label{tab:baseline-ablation}
\setlength{\tabcolsep}{4pt}
\begin{tabularx}{\linewidth}{lYrrr}
\toprule
对照层 & 计划与配置 & 平均工期/cycle & A0速度比均值 & 配对效应比均值\\\midrule
统一基线 & 整图单Task、单核$A0$ & 3124793.73 & 1.0000 & --\\
问题一终态 & $A(X_A^{\mathrm{inh}})$ & 921892.21 & 3.4132 & --\\
问题二终态 & $B(X_B^{\mathrm{inh}})$ & 903275.96 & 3.3862 & --\\
问题三终态 & $C(X_C^{\mathrm{inh}})$ & 895871.90 & 3.4272 & --\\
同计划无L2 & 冻结配对中的$B(X_B^*)$ & 929324.47 & -- & 1.0000\\
固定计划开启L2 & $C(X_B^*)$ & 922208.33 & -- & 1.0125\\
固定L2调整计划 & $C(X_C^*)$ & 922120.10 & -- & 1.0011\\
\bottomrule
\end{tabularx}
\end{table}
平均工期先对100张图的完成周期取算术平均；两个比值列先逐图计算，再对100个比值取算术平均。前三行采用逐图继承后的终态，后四行保留固定五核计划的受控配对。最后一列中，1.0125是固定计划的$F_B(X_B^*)/F_C(X_B^*)$均值，1.0011是固定C配置的$F_C(X_B^*)/F_C(X_C^*)$均值；其分子、分母与A0速度比列不同。固定计划开启L2使54图工期缩短、46图持平；固定C配置再调度使7图进一步改善。两步合成的比值均值为1.0136，逐图定义见式\eqref{eq:cache-factors}。

候选评价的阶段结果见表~\ref{tab:candidate-status}。表中“组合数”是最终结果的1500个图--核数--场景组合；其余列的分母是实际生成并记录的候选次数，同一组合可以贡献多条候选记录。A场景有214条候选在profile阶段失败，B、C场景没有单独记录的profile失败；B和C分别有110条和77条候选在完整事件评价阶段失败。成功数与失败数相加得到各场景的候选尝试总数，1500个最终组合均保留一条成功的官方结果。
\begin{table}[H]\centering\small
\caption{主矩阵候选评价的阶段分层与分母}\label{tab:candidate-status}
\setlength{\tabcolsep}{4pt}
\begin{tabular}{lrrrrr}\toprule
场景 & 组合数 & 官方结果成功候选 & profile失败 & 完整评价失败 & 候选尝试数\\\midrule
A & 500 & 1000 & 214 & 0 & 1214\\
B & 500 & 890 & 0 & 110 & 1000\\
C & 500 & 894 & 0 & 77 & 971\\\midrule
合计 & 1500 & 2784 & 214 & 187 & 3185\\\bottomrule
\end{tabular}
\end{table}
这里的“官方结果成功”表示候选通过前置检查并由题给事件评价器返回完整结果；“完整评价失败”表示候选已经进入事件评价阶段但未返回工期。2784/3185是候选记录的成功占比，最终组合覆盖率另为1500/1500。

\subsubsection{问题一结果}
问题一模型把张量的消费Task集合写入结构搬运公式，并把同核、跨核激活等待送入时间递推。按逐图继承规则，五核平均加速比为3.4132，中位数3.9605，100张图均不低于整图单核基准。均值和中位数同时报告，是因为两者回答不同问题：前者把100个速度比等权平均，后者描述第50个左右的图；两项共同展示总体水平与图间离散性，单图结论见逐图结果。

以每个图--核数--场景的首个成功候选为起点，问题一正式400组中117组在追加候选评价后得到更短的最终方案，平均相对工期降幅为10.5562\%。该差值是候选搜索的选择收益；它既不是相对统一A0基线的方案收益，也不是移除某个机制后的消融。五核单独看，39图改善，平均工期降幅16.9668\%，对应的是五核子集的搜索收益。

原始候选的工期可能随核数增加而上升，逐图继承后则不再出现这种退化。\Case{067}的A场景四核工期为18473155 cycle，原始五核候选为63959635 cycle；最终五核沿用四核计划并补齐空核列表，因而该图被纳入不退化的终态统计。

\subsubsection{问题二结果}
问题二把同核子图合并成核心级驻留域，结构COPY按远端接收核心去重，同时仍受L1和UB物理版本生存区间制约。按逐图继承规则，五核平均加速比为3.3862，中位数4.0046，100图均不低于整图单核。固定五核计划的A/B对照另在受控消融中报告。

上述逐图关系已在问题二的图\ref{fig:b-ab-tradeoff}中给出，本节沿用同一100图配对结果，不再重复排图。

对\Case{053}，B最终方案的额外搬运为5405536 B，且全部记录为换出恢复搬运；同核长期驻留降低了某些结构边界，但容量恢复和执行顺序仍可能扩大工期。该例同时包含驻留规则和计划变化，规则效应由固定计划的受控对照给出。

在固定核数的原始候选比较中，问题二500组有83组由首个成功候选到最终候选得到更短工期，平均搜索收益为2.7576\%；五核子集有18图改善，平均降幅6.7385\%。该比较只描述候选搜索收益，与统一A0基线及受控机制对照分开报告。模型下界拆为计算下界$LB_{\rm comp}$、I/O下界$LB_{\rm io}$，并取$LB=\max\{LB_{\rm comp},LB_{\rm io}\}$；前者由依赖最长路与M、V流水线工作量给出，后者由必需I/O工作量与DDR带宽给出。两项均未包含完整事件中的等待、容量换出和共享资源争用。表\ref{tab:eval-cost}的$F/LB$统一采用聚合下界$LB$，正式B的500个最终工期与该下界的比值中位数为1.3491；该下界用于排查距离较大的方案，最终选优仍按事件评价工期比较。

\subsubsection{问题三结果}
问题三用相同B最终计划分别在B/C评价，再在C中选择最终计划，因而式\eqref{eq:cache-factors}的三段比值具有清楚的身份条件。500个配对中，硬件效应改善205组、持平295组；C内调度适配改善47组、持平453组；综合改善232组、持平268组，三个集合的中位比值均为1。

五核固定计划下54图硬件改善、46图持平、零图退化，$C(X_B^*)$字节加权命中率为16.79\%。54张改善图的平均工期降幅为2.0074\%，对应的速度增幅均值为2.31\%；两项分别对应速度比与工期降幅，按各自定义解读。\Case{080}的硬件效应达到1.2549，固定计划工期降幅为20.3138\%，而\Case{055}虽有189952字节命中，完成时间没有变化。命中只是服务来源，实际是否减少尾部完成时间取决于事件位置和资源竞争。五核C内进一步适配仅改善7图，平均适配比1.0011；这与轻量代理的用途一致：代理只负责候选排序，精确FIFO查询、插入、淘汰和工期仍由官方事件模拟决定，未改善案例需结合官方事件字段解释，代理只决定候选试评顺序。

问题三正式500组中47组由首个成功候选到最终候选得到更短工期，平均搜索收益为0.5006\%；五核7图改善，平均降幅0.1098\%。该指标只表示有限候选中的选优收益，机制作用仍由表\ref{tab:baseline-ablation}中的固定计划对照衡量。

\subsubsection{依赖结构分组结果}
全量均值可能掩盖不同计算图的调度难度。以第4章从原图得到的弱连通分量数分层，再对每组的五核最终方案计算速度比均值及低于1的图数，结果见表\ref{tab:topology-speedup}。分组只使用原始输入，不根据求解结果挑选案例；每个场景的三组样本数之和均为100。
\begin{table}[H]\centering\small
\caption{原始计算依赖分量数与五核最终方案速度比的分组关系}\label{tab:topology-speedup}
\setlength{\tabcolsep}{4pt}
\begin{tabular}{lrrrrrrr}\toprule
原始弱连通分量 & 图数 & A均值 & B均值 & C均值 & A低于1 & B低于1 & C低于1\\\midrule
1个 & 16 & 1.2963 & 1.1660 & 1.1757 & 0 & 0 & 0\\
2--100个 & 57 & 3.3607 & 3.3496 & 3.3982 & 0 & 0 & 0\\
超过100个 & 27 & 4.7783 & 4.7791 & 4.8227 & 0 & 0 & 0\\\bottomrule
\end{tabular}\end{table}

单分量图仍需沿依赖切开大分量，多分量图则有更多可直接装箱的独立块；分量数与规模、依赖形状共同影响收益。逐图继承规则把整图基准纳入每个核数的候选，因此三类拓扑分组的最终低速图数均为0；表中均值沿用各组继承后的速度比口径，具体差异仍由官方事件时间线决定。

\subsubsection{低速图的逐例分析}
分组均值之外，还需检查同一张图在三问中是否跨过单核基准。以五核最终工期的真实速度比$S<1$记为“慢”，逐图匹配100个案例，得到表\ref{tab:slow-transition}。四种状态覆盖全部100图；其中C与B的慢例集合相同，C场景的平均速度比上升与这些难例的逐图状态分别解读。
\begin{table}[H]\centering\small
\caption{五核100图在A、B、C最终方案中的低速状态交叉}\label{tab:slow-transition}
\begin{tabular}{lllr}\toprule
A & B & C & 图数\\\midrule
不低于单核 & 不低于单核 & 不低于单核 & 100\\\bottomrule
\end{tabular}\end{table}

逐图继承后，A、B、C三场景五核终态均不低于单核基准，表\ref{tab:slow-transition}只保留一个状态。固定计划下的场景差异仍通过表\ref{tab:baseline-ablation}和第8章的同计划配对解释；继承规则负责消除已评价候选中的工期退化，不改变候选生成或官方事件评价。

\subsection{模型的不足}
候选搜索为每个图--核数--场景保留有限的完整评价机会；候选构造与轻量排序决定了实际搜索范围，大连通分量、容量临界点和重复Cache失效可能需要更多候选。主矩阵的分层记录见表\ref{tab:candidate-status}；官方单核基准覆盖100/100图，B/C同图同核配对500组，结果字段检查、15项抽样复算及3项程序单元测试均通过。代表图的关键瓶颈由事件时间线逐例解释。

表\ref{tab:eval-cost}从最终选中方案的结果字段统计聚合下界距离和单次官方评价耗时，每个场景500组，问题一包含用于诊断的100个单核组合。表中$LB=\max\{LB_{\rm comp},LB_{\rm io}\}$；$LB_{\rm comp}$保留依赖最长路与M、V两条计算流水线的工作量，$LB_{\rm io}$保留必需I/O工作量与DDR带宽约束，两项均不包含完整事件中的同步等待、容量换出和共享资源争用。$F/LB$的上十分位数超过4，该比值用于发现离下界较远的图，算法选优仍由官方结果完成。三个场景的中位距离只能在各自执行规则内比较。
\begin{table}[H]\centering\small
\caption{已选方案的参考下界距离与单次官方评价耗时}\label{tab:eval-cost}
\setlength{\tabcolsep}{4pt}
\begin{tabular}{lrrrrr}\toprule
场景 & 组合数 & \shortstack{$F/LB$\\中位数} & \shortstack{$F/LB$\\第90百分位} & \shortstack{评价耗时\\中位数/s} & \shortstack{评价耗时\\第90百分位/s}\\\midrule
A & 500 & 1.325 & 4.617 & 1.484 & 53.487\\
B & 500 & 1.349 & 4.486 & 0.755 & 6.613\\
C & 500 & 1.314 & 4.486 & 0.752 & 6.618\\\bottomrule
\end{tabular}\end{table}

表中秒数是最终方案的一次官方评价耗时，不包括候选构造与前序评价。A场景第90百分位为53.487~s、中位数为1.484~s，显示大图评价成本与典型图相差较大；当前流程按必评锚点和有限调用额度安排完整评价，轻量排序用于确定其余候选的试评顺序。当前计时没有覆盖从输入读取、候选生成到最终写出的统一端到端流程，表中秒数只说明候选评价成本；题面建议时限需要端到端计时。

\subsection{候选级算法效率与计算成本}
算法成本分为图索引、候选构造与评分、官方评价三个阶段。令$|V|$为非COPY操作数，$|E_0|$为原始边数，$|E|$为去重后的计算依赖数，$|T|$为张量数，$P$为逐张量枚举的生产者--消费者配对数。生产消费索引遍历原始边并展开这些配对，代价为$O(|V|+|T|+|E_0|+P)$；堆拓扑序及其后继排序为$O((|V|+|E|)\log |V|)$。弱连通分量的邻接遍历为$O(|V|+|E|)$，但当前实现逐分量调用\texttt{min(left)}选起点，最坏另需$O(|V|^2)$。释放优先序每步对全部就绪操作排序后取前8项，最坏为$O(|V|^2\log |V|)$，不能归入线性预处理。

若生成$M$个候选，覆盖与依赖检查、计划去重和轻量评分逐候选执行；评分还扫描张量接收关系并排序容量区间端点。记$R$为张量生产消费关联数，$K$为核数，单候选成本为$C_{\rm cand}(|V|,|E|,|T|,R,K)$，则候选阶段为$O(MC_{\rm cand}+M\log M)$。官方profile与事件模拟另行展开Task、COPY和物理版本。令$E_{\rm ev}$为展开后的事件数，即使事件队列操作约为$O(E_{\rm ev}\log E_{\rm ev})$，Step2/3容量扫描和服务池状态处理仍需单独计入，不能由队列复杂度代表整次评价。表\ref{tab:algorithm-timing}报告有计时字段的候选实测成本；每个图--核数--场景最多进入两次完整官方评价。
\begin{table}[H]\centering\small
\caption{有完整计时字段的候选的官方评价耗时（秒）}
\label{tab:algorithm-timing}
\setlength{\tabcolsep}{2.4pt}
\begin{tabular}{lrrrrr}\toprule
场景 & 成功候选数 & \shortstack{profile\\中位数} & \shortstack{profile\\P95} & \shortstack{完整评价\\中位数} & \shortstack{完整评价\\P95/最大值}\\\midrule
A & 302 & 1.196 & 5.310 & 3.011 & 93.317 / 384.720\\
B & 275 & 1.578 & 8.684 & 1.680 & 9.315 / 22.891\\
C & 280 & 1.534 & 8.753 & 1.764 & 8.776 / 15.102\\\bottomrule
\end{tabular}
\end{table}
表\ref{tab:algorithm-timing}统计结果文件中有完整计时字段的成功候选，P95和最大值按候选记录计算；A场景的大图尾部明显更长，说明按图规模和展开事件数分配评价预算比固定平均耗时更合适。\path{timing.csv}共记录963条候选：A场景369条中302条成功、67条profile失败，B场景302条中275条成功、27条完整评价失败，C场景292条中280条成功、12条完整评价失败。这963条构成独立的计时子集，与表~\ref{tab:candidate-status}的主矩阵分层分别统计。进一步按\path{candidate_evaluations.csv}将同一图--核数--场景的候选计时相加，A/B/C的组合级中位数分别为3.9917/2.4187/2.3171~s，P95为225.4680/51.5057/39.8999~s，最大值为1198.4254/211.7056/212.4157~s；三场景候选计时总和分别为21759.58/4845.95/4268.67~s。上述计时覆盖profile级、完整官方评价级和组合级候选串行计时，输入读取、候选生成、前序评价及最终写出未计入，因此本文将其作为候选级效率证据；轻量排序把昂贵的官方事件模拟限制在每组有限候选内，效率收益来自减少完整评价次数。

其中，问题一出现了约1198~s（约20~min）的组合级候选评价耗时，超过题面脚注给出的5--10~min建议。该数值由候选评价耗时相加，整场程序从输入到输出的端到端墙钟尚无对应计时；它同时显示大图评价成本的尾部。对这类超大图，实际运行应先按非COPY操作数、展开事件数和内存压力设定候选数量上限，保留整图基准、首个合法分块和当前最优方案；若前两类候选的下界距离已稳定，则停止追加完整评价，并把未评候选记入未评台账。本文当前结果仍按统一候选预算报告，超时案例作为效率边界保留。

\subsection{模型改进与推广}
模型适用于依赖图、操作成本及片上容量已知，且能够对完整计划作事件评价的调度任务。结构COPY式刻画接收域，物理版本生存区间对应容量换出，FIFO查询与完成事件决定共享读服务。对本题而言，可沿三个方向继续推进：增加容量临界与共享输入时序候选，建立能回放FIFO事件的轻量评分，并对胜出与失效的计划逐事件定位关键路径。机制移除对照可进一步区分各项改动的作用；跨设备性能则需相应平台的实际测量。本文给出的数值结论对应100图、题设硬件配置和官方CPU Python事件模拟下的已评价候选集合。
